\documentclass[11pt]{article}

\usepackage[T1]{fontenc}
\usepackage{lmodern}
\usepackage[margin=1in]{geometry}
\usepackage{amsmath}
\usepackage{amssymb}
\usepackage{booktabs}
\usepackage{longtable}
\usepackage{array}
\usepackage{listings}
\usepackage{xcolor}
\usepackage{hyperref}

\hypersetup{
  colorlinks=true,
  linkcolor=blue!50!black,
  citecolor=blue!50!black,
  urlcolor=blue!60!black
}

\lstset{
  basicstyle=\ttfamily\footnotesize,
  columns=fullflexible,
  breaklines=true,
  frame=single,
  framesep=4pt,
  xleftmargin=2pt,
  xrightmargin=2pt,
  keywordstyle=\color{blue!60!black}\bfseries,
  commentstyle=\color{gray!70!black}\itshape,
  stringstyle=\color{green!50!black},
  showstringspaces=false
}

\title{The Text Property Graph: A Code-Inspired Representation for Security Text}
\author{}
\date{}

\begin{document}
\maketitle

\section{Introduction}

Security text is not only a sequence of words. A CVE description, exploit
report, vendor advisory, or threat-intelligence note usually states which
vulnerability is being discussed, which software is affected, which version is
vulnerable, how the attack works, what impact follows, and what mitigation is
recommended. These facts are not isolated keywords. They are connected through
syntax, sentence order, entity references, discourse markers, and
domain-specific security relations.

The Text Property Graph (TPG) is the representation used in this project to
preserve those relations. A TPG is a directed, typed, attributed property
graph. Its nodes represent linguistic and security objects, such as documents,
paragraphs, sentences, tokens, entities, predicates, CVE identifiers, software
products, versions, attack vectors, impacts, severities, and remediations. Its
edges represent typed relations between those objects, including containment,
dependency, sequence, coreference, predicate-argument binding, discourse
relation, entity relation, and security-specific relations.

The TPG is inspired by the Code Property Graph (CPG). A CPG combines several
views of source code into one graph: syntax, control flow, data flow, and
dependence. The TPG transfers the same design principle to text. Instead of
leaving dependency parses, named entities, coreference links, discourse
relations, and security facts as disconnected outputs, it stores them as
different edge types over one shared graph.

This chapter explains the problem that motivates TPG, the fundamental structure
of the representation, its relationship to the CPG, and why it is useful for
security-related text. The description is grounded in the project
documentation and implementation, especially the files under
\texttt{docs/}, \texttt{EPSS\_TPG/docs/}, and \texttt{EPSS\_TPG/tpg/}.

\section{Problem Definition}

\subsection{Why Flat Text Representations Are Not Enough}

Most modern NLP pipelines begin by reducing text to tokens or subword tokens.
Those tokens are then embedded by a transformer and pooled or passed to a
downstream model. This approach is effective for many tasks, but it does not
make the structure of the text explicit. A vector representation may contain
useful information, but it does not directly expose relations such as:

\begin{itemize}
    \item which CVE affects which software product;
    \item which version belongs to that software;
    \item which attack vector leads to which impact;
    \item which sentence gives the mitigation;
    \item which pronoun or phrase refers back to a previous entity;
    \item which sentence explains, contrasts with, or follows another sentence.
\end{itemize}

For security text, this is a serious limitation. A vulnerability advisory is
often short and dense. Its meaning depends on relations. The words
\texttt{remote}, \texttt{critical}, \texttt{Apache}, and \texttt{patch} are
useful, but the important question is how they are connected. A sentence saying
that a remote attacker can exploit Apache is different from a sentence saying
that Apache released a patch to prevent remote exploitation. A representation
that only stores a token sequence has to infer this structure indirectly.

\subsection{What the TPG Is Designed to Preserve}

The TPG is designed to preserve the structural information that a flat token
sequence hides. It keeps the following layers in one graph:

\begin{itemize}
    \item document structure, from document to paragraph to sentence to token;
    \item syntactic dependency structure inside each sentence;
    \item sequential flow across tokens, sentences, and paragraphs;
    \item entity mentions and coreference chains;
    \item predicate-argument relations;
    \item discourse and rhetorical relations;
    \item security entities and security-specific typed relations;
    \item contextual embeddings when the model-based frontend is available.
\end{itemize}

The goal is not to replace transformer embeddings. The goal is to place those
embeddings inside a richer structure. In the EPSS-TPG pipeline, the final graph
can be exported to PyTorch Geometric as \texttt{x}, \texttt{edge\_index},
\texttt{edge\_type}, and \texttt{edge\_attr}, so a GNN can use both node
features and typed graph structure.

\section{Why We Use a Property Graph}

\subsection{Multiple Relations on the Same Text}

Text contains many relation types at the same time. A token belongs to a
sentence. The same token may depend syntactically on another token, appear
inside an entity span, be followed by another token, and carry a contextual
embedding. A sentence may continue the previous sentence, contrast with it, or
explain it. A CVE mention may be connected to a software product, a CWE, an
impact, and a mitigation.

A property graph is a natural fit because it allows all of these relations to
coexist. The same node can participate in several edge types without being
duplicated. This is the main reason TPG is more expressive than a single
sequence, a single parse tree, or a bag of extracted entities.

\subsection{Typed Nodes, Typed Edges, and Properties}

In a TPG, every node has a type and a property map. The implementation stores
node properties such as text, lemma, part-of-speech tag, dependency relation,
entity type, sentence index, token index, character offsets, domain type,
confidence, source, and optional embedding. Every edge also has a type and may
carry properties such as dependency label, semantic-role label, discourse
label, coreference cluster, relation type, confidence, or other metadata.

This follows the same property-graph idea used by Joern's CPG. The project
implementation stores this schema in \texttt{tpg/schema/types.py}, with
\texttt{NodeProperties}, \texttt{EdgeProperties}, \texttt{NodeType},
\texttt{EdgeType}, \texttt{SecurityNodeType}, \texttt{SecurityEdgeType}, and
\texttt{TPGSchema}.

\section{spaCy as the Linguistic Backend}

\subsection{Official Definition and Role}

The generic TPG does not start from hand-written token rules. It starts from
spaCy, which is the linguistic backend of the Level 1 graph construction.
spaCy describes itself as ``Industrial-strength Natural Language Processing in
Python''\footnote{\url{https://spacy.io/}}. Its project description defines it
as a library for advanced NLP in Python and Cython, designed for real products,
with pretrained pipelines for tokenization, tagging, parsing, named entity
recognition, text classification, and related NLP tasks.\footnote{\url{https://pypi.org/project/spacy/}}

The important point for TPG is that spaCy does not only split text into tokens.
Its standard processing pipeline converts raw text into a \texttt{Doc} object
with linguistic annotations. The official linguistic-features documentation
states that spaCy is designed so that raw text becomes a \texttt{Doc} object
with annotations attached to it.\footnote{\url{https://spacy.io/usage/linguistic-features\#\#tokenization}}
Those annotations include token boundaries, sentence boundaries,
part-of-speech tags, lemmas, dependency-parse labels, noun chunks, and named
entities when the selected pipeline provides the relevant components.

\subsection{How spaCy Works}

At a high level, spaCy processes text through a language pipeline. The first
stage is tokenization. spaCy's tokenizer is non-destructive: the tokenized
\texttt{Doc} can reconstruct the original input text. This matters for TPG
because token character offsets must remain aligned with the original document.

After tokenization, trained spaCy components can add statistical annotations.
For English pipelines such as \texttt{en\_core\_web\_sm}, the relevant
annotations are:

\begin{itemize}
    \item \textbf{sentence boundaries}, exposed through \texttt{Doc.sents};
    \item \textbf{tokens}, with text, lemma, part-of-speech tag, dependency
    relation, token index, and character offsets;
    \item \textbf{dependency parses}, where each token has a syntactic head and
    a dependency label such as \texttt{nsubj}, \texttt{dobj},
    \texttt{amod}, or \texttt{prep};
    \item \textbf{named entities}, exposed through \texttt{Doc.ents}, with
    entity labels such as organization, product, date, or geopolitical entity;
    \item \textbf{noun chunks}, which spaCy defines as base noun phrases
    available through \texttt{Doc.noun\_chunks} when a parser is present.
\end{itemize}

The dependency parser is especially important for TPG. spaCy's dependency
parse gives the syntactic backbone that becomes the \texttt{DEP} subgraph. In
the CPG analogy, this is the closest text-side equivalent of the AST view.

\subsection{How TPG Uses spaCy}

The project uses spaCy through \texttt{tpg/frontends/spacy\_frontend.py}. The
\texttt{SpacyFrontend} attempts to load \texttt{en\_core\_web\_sm}. If that
model is unavailable, it falls back to a blank English pipeline with a
sentencizer. The fallback can still split sentences, but it does not provide
the same parser, POS tagger, lemmatizer, noun chunks, or named entity
recognition as the trained model. Therefore, the full TPG linguistic structure
depends on the trained spaCy pipeline being available.

The TPG construction uses spaCy annotations as follows:

\begin{enumerate}
    \item The input text is split into paragraphs using blank-line boundaries.
    \item Each paragraph is passed to spaCy, producing a \texttt{Doc}.
    \item For every sentence in \texttt{Doc.sents}, the frontend creates a
    \texttt{SENTENCE} node.
    \item For every spaCy token, it creates a \texttt{TOKEN} node with text,
    lemma, POS tag, dependency relation, entity IOB tag, entity type, sentence
    index, token index, and character span.
    \item For every named entity in \texttt{sent.ents}, it creates an
    \texttt{ENTITY} node and links that entity to its tokens with
    \texttt{BELONGS\_TO}.
    \item For content verbs, it creates \texttt{PREDICATE} nodes. Auxiliary
    verbs are filtered out.
    \item For syntactic children whose dependency labels can be mapped to
    semantic roles, it creates \texttt{ARGUMENT} nodes and \texttt{SRL\_ARG}
    edges.
    \item For noun chunks, it creates \texttt{NOUN\_PHRASE} nodes.
    \item For selected clausal dependencies such as \texttt{advcl},
    \texttt{relcl}, \texttt{ccomp}, and \texttt{acl}, it creates
    \texttt{CLAUSE} nodes.
    \item For verb phrases with auxiliaries, negation, particles, or adverbial
    modifiers, it creates \texttt{VERB\_PHRASE} nodes.
\end{enumerate}

The spaCy frontend also creates the initial graph edges. It adds
\texttt{CONTAINS} edges for document hierarchy, \texttt{DEP} edges from each
token's syntactic head to its dependent, \texttt{NEXT\_TOKEN} edges for token
order, \texttt{NEXT\_SENT} edges for sentence order, \texttt{NEXT\_PARA} edges
for paragraph order, and cross-sentence \texttt{NEXT\_TOKEN} edges from the
last token of one sentence to the first token of the next sentence.

\subsection{Why spaCy Is the Right Backend for TPG}

spaCy is suitable for this role because it gives the TPG a stable linguistic
front end. It provides tokenization, sentence segmentation, POS tags,
dependency labels, entity spans, noun chunks, lemmas, and character offsets in
one consistent \texttt{Doc} object. The TPG frontend then converts those
annotations into graph nodes and edges.

In other words, spaCy supplies the first layer of linguistic analysis, while
TPG supplies the graph representation. spaCy answers: what are the tokens,
sentences, syntactic heads, noun chunks, and entities? TPG answers: how should
those objects be connected in a typed graph that can be queried, exported, and
used by a GNN?

Once the spaCy frontend has produced the linguistic skeleton, a separate
embedding pass attaches contextual vectors. The project uses
\textbf{SecBERT}\footnote{\url{https://huggingface.co/jackaduma/SecBERT}}, a
BERT model further pre-trained on cybersecurity text, kept frozen at inference
time. The embedder runs SecBERT on the raw text, aligns its WordPiece
sub-tokens to spaCy token character spans, and stores a 768-dimensional
contextual vector on each \texttt{TOKEN} node under
\texttt{node.properties.extra["embedding"]}. Sentence-level
\texttt{[CLS]}-derived embeddings and entity-level mean embeddings are stored
on \texttt{SENTENCE} and \texttt{ENTITY} nodes by the same pass. Downstream,
the GNN reads each node's feature as the concatenation of a 13-dimensional
node-type one-hot and that 768-dimensional embedding, giving a 781-dimensional
node-feature vector. The detailed architecture of the GNN that consumes these
features is described in the companion document
\texttt{multiview\_tpg\_formula.tex}.

\section{Fundamental Concepts of the TPG}

\subsection{Formal Definition}

A Text Property Graph for one document can be written as:

\[
G = (V, E, X, R, \tau, \pi_V, \pi_E)
\]

where \(V\) is the set of nodes, \(E\) is the set of directed edges, \(X\) is
the node-feature matrix, \(R\) is the edge-relation vocabulary, \(\tau\) maps
each edge to a relation type, and \(\pi_V\), \(\pi_E\) are node and edge
property maps.

In the PyTorch Geometric export, the graph is represented by:

\begin{itemize}
    \item \texttt{x}: node features, including node-type one-hot features and
    optional contextual embeddings;
    \item \texttt{edge\_index}: directed graph connectivity;
    \item \texttt{edge\_type}: integer edge-type labels;
    \item \texttt{edge\_attr}: one-hot edge-type features;
    \item \texttt{y}: graph-level label when the graph is used for supervised
    learning.
\end{itemize}

\subsection{Two Practical Schema Levels}

For this project, the TPG is used at two practical levels.

\begin{table}[h]
\centering
\small
\begin{tabular}{lp{0.68\textwidth}}
\toprule
Level & Description \\
\midrule
Level 1: Generic TPG & General text graph with 13 base node types and 13 base edge types. It is produced mainly by the spaCy frontend and enrichment passes. \\
Level 2: Security TPG & Security-aware graph that adds 12 security node types and 10 security edge types for vulnerability text. \\
\bottomrule
\end{tabular}
\caption{TPG levels used in this chapter and in the EPSS-TPG work.}
\end{table}

The generic level is meant for ordinary text. The security level is the
domain-specific extension used for CVE advisories and vulnerability-related
text. This chapter focuses only on these two levels.

\section{Connection to the Code Property Graph}

\subsection{What the CPG Solved for Code}

The Code Property Graph was introduced for source-code analysis. A code
vulnerability often cannot be detected from syntax alone. One may need to know
where a function call appears, whether a guard exists before it, whether
untrusted input reaches it, and which statements depend on which conditions.
The CPG solves this by storing several views of the program in one property
graph.

The TPG applies the same idea to text. The analogy is not that natural language
is the same as source code. The analogy is that both code and text contain
several structural views, and these views are more useful when they are joined
into one graph.

\subsection{CPG-to-TPG Mapping}

The mapping from CPG concepts to TPG concepts is implemented directly in the
project schema. The main mapping is shown below.

\begin{table}[h]
\centering
\small
\begin{tabular}{lll}
\toprule
CPG view or edge & TPG analogue & Meaning in text \\
\midrule
AST & \texttt{DEP} & Syntactic dependency structure \\
CFG (intra-block) & \texttt{NEXT\_TOKEN} & Token order inside a sentence \\
CFG (cross-block) & \texttt{NEXT\_SENT} & Sentence order within a paragraph \\
CFG (cross-scope) & \texttt{NEXT\_PARA} & Paragraph order within the document \\
REACHING\_DEF & \texttt{COREF} & Entity reference flow through text \\
ARGUMENT & \texttt{SRL\_ARG} & Predicate-argument binding \\
CDG & \texttt{RST\_RELATION} & Rhetorical relation: cause, contrast, condition \\
DOMINATE & \texttt{DISCOURSE} & General discourse continuation \\
CONTAINS & \texttt{CONTAINS} & Structural containment \\
BINDS\_TO & \texttt{BELONGS\_TO} & Token membership in entity or phrase \\
CALL & \texttt{ENTITY\_REL} & Relation between entities through predicates \\
EVAL\_TYPE & \texttt{SIMILARITY} & Semantic similarity between nodes \\
(no CPG analogue) & \texttt{AMR\_EDGE} & AMR semantic relation slot \\
\bottomrule
\end{tabular}
\caption{Main CPG-to-TPG mapping used in the implementation. Source:
\texttt{tpg/schema/types.py}, dictionaries \texttt{JOERN\_TO\_TPG\_NODE\_MAP}
and \texttt{JOERN\_TO\_TPG\_EDGE\_MAP}.}
\end{table}

This mapping is the core reason TPG can be described as code-inspired. It is
not just a graph built from text; it is a graph whose views deliberately mirror
the multi-view structure of CPGs.

\section{Nodes}

\subsection{Generic Node Types}

The generic TPG schema defines 13 node types.

\begin{longtable}{lp{0.68\textwidth}}
\caption{Generic TPG node types.}\\
\toprule
Node type & Role \\
\midrule
\endfirsthead
\toprule
Node type & Role \\
\midrule
\endhead
\bottomrule
\endfoot
\texttt{DOCUMENT} & Root node for the whole document. \\
\texttt{PARAGRAPH} & Structural container for a paragraph. \\
\texttt{SENTENCE} & Sentence-level unit. \\
\texttt{TOKEN} & Atomic linguistic token. \\
\texttt{ENTITY} & Named or extracted entity. \\
\texttt{PREDICATE} & Verb or action node. \\
\texttt{ARGUMENT} & Role-bearing argument linked to a predicate. \\
\texttt{CONCEPT} & Abstract concept or category. \\
\texttt{NOUN\_PHRASE} & Noun phrase chunk. \\
\texttt{VERB\_PHRASE} & Verb phrase chunk. \\
\texttt{CLAUSE} & Clause or subordinate structure. \\
\texttt{MENTION} & Coreference mention such as a pronoun. \\
\texttt{TOPIC} & Document-level topic node. \\
\end{longtable}

\subsection{Security Node Types}

The security extension adds 12 domain-specific node types.

\begin{longtable}{lp{0.68\textwidth}}
\caption{Security-specific TPG node types.}\\
\toprule
Node type & Role \\
\midrule
\endfirsthead
\toprule
Node type & Role \\
\midrule
\endhead
\bottomrule
\endfoot
\texttt{CVE\_ID} & Vulnerability identifier. \\
\texttt{CWE\_ID} & Weakness class identifier. \\
\texttt{SOFTWARE} & Affected software product. \\
\texttt{VERSION} & Software version. \\
\texttt{COMPONENT} & Affected component or module. \\
\texttt{CODE\_ELEMENT} & Function, API, variable, or code construct mentioned in text. \\
\texttt{ATTACK\_VECTOR} & Attack method or access path. \\
\texttt{IMPACT} & Consequence of exploitation. \\
\texttt{REMEDIATION} & Patch, upgrade, workaround, or mitigation. \\
\texttt{THREAT\_ACTOR} & Attacker or actor type. \\
\texttt{VULN\_TYPE} & Vulnerability category such as buffer overflow, XSS, or SQL injection. \\
\texttt{SEVERITY} & CVSS score or qualitative severity term. \\
\end{longtable}

\section{Edges}

\subsection{Generic Edge Types}

The edge vocabulary is the most important part of the TPG. It is what turns a
set of text objects into a multi-view graph.

\begin{longtable}{lp{0.68\textwidth}}
\caption{Generic TPG edge types.}\\
\toprule
Edge type & Meaning \\
\midrule
\endfirsthead
\toprule
Edge type & Meaning \\
\midrule
\endhead
\bottomrule
\endfoot
\texttt{DEP} & Syntactic dependency from head token to dependent token. \\
\texttt{NEXT\_TOKEN} & Sequential token order. \\
\texttt{NEXT\_SENT} & Sequential sentence order. \\
\texttt{NEXT\_PARA} & Sequential paragraph order. \\
\texttt{COREF} & Coreference link between repeated mentions or pronouns and entities. \\
\texttt{SRL\_ARG} & Predicate-argument relation. \\
\texttt{AMR\_EDGE} & Abstract meaning relation slot in the schema. \\
\texttt{RST\_RELATION} & Explicit rhetorical relation such as cause, contrast, temporal, condition, or summary. \\
\texttt{DISCOURSE} & General discourse continuation or topic-to-sentence relation. \\
\texttt{CONTAINS} & Structural containment from document to paragraph to sentence to node. \\
\texttt{BELONGS\_TO} & Membership relation from entity, phrase, predicate, or argument to token. \\
\texttt{ENTITY\_REL} & Relation between entities, usually inferred from nearby predicates. \\
\texttt{SIMILARITY} & Semantic similarity relation. \\
\end{longtable}

\subsection{Security Edge Types}

Security edges are optional first-class edge types emitted by
\texttt{SecurityRelationsPass}. They are not only properties on generic
\texttt{ENTITY\_REL} edges.

\begin{table}[h]
\centering
\small
\begin{tabular}{lll}
\toprule
Security edge & Source role & Target role \\
\midrule
\texttt{SEC\_AFFECTS} & CVE / vulnerability & software \\
\texttt{SEC\_HAS\_VERSION} & software & version \\
\texttt{SEC\_LOCATED\_IN} & CVE / vulnerability & component \\
\texttt{SEC\_CLASSIFIED\_AS} & CVE & CWE \\
\texttt{SEC\_EXPLOITED\_BY} & CVE / vulnerability & attack vector \\
\texttt{SEC\_CAUSES} & attack vector & impact \\
\texttt{SEC\_MITIGATED\_BY} & CVE / vulnerability & remediation \\
\texttt{SEC\_USES\_FUNCTION} & CVE / vulnerability & code element \\
\texttt{SEC\_THREATENS} & attack vector / threat actor & software \\
\texttt{SEC\_HAS\_SEVERITY} & CVE & severity \\
\bottomrule
\end{tabular}
\caption{Security relation types used by the Level 2 TPG.}
\end{table}

When security edges are enabled, the edge vocabulary expands from 13 generic
edge types to 23 total edge types. Security edges occupy indices 13--22 in the
unified edge-type vocabulary. This matters for GNN training: the model can only
learn from these relations when the dataset and model use the expanded
edge-type vocabulary.

\section{How the TPG Is Constructed}

\subsection{Generic Pipeline}

The generic pipeline is implemented by \texttt{TPGPipeline}. It follows the
same broad shape as Joern: a frontend builds the initial graph, enrichment
passes add more edges and nodes, and exporters write the graph to external
formats.

\begin{enumerate}
    \item The spaCy frontend parses the text and creates document, paragraph,
    sentence, token, entity, predicate, argument, phrase, and clause nodes.
    \item The frontend adds structural and syntactic edges such as
    \texttt{CONTAINS}, \texttt{DEP}, \texttt{NEXT\_TOKEN},
    \texttt{NEXT\_SENT}, \texttt{NEXT\_PARA}, \texttt{BELONGS\_TO}, and
    \texttt{SRL\_ARG}.
    \item \texttt{CoreferencePass} adds \texttt{COREF} edges.
    \item \texttt{DiscoursePass} adds \texttt{RST\_RELATION} or
    \texttt{DISCOURSE} edges.
    \item \texttt{EntityRelationPass} adds \texttt{ENTITY\_REL} edges between
    entities in the same sentence, using nearby predicates when available.
    \item \texttt{TopicPass} creates \texttt{TOPIC} nodes and links them to
    relevant sentences.
    \item The SecBERT embedder runs once over the raw text, aligns its
    WordPiece sub-tokens to spaCy token character spans, and attaches a
    768-dimensional contextual vector to every \texttt{TOKEN} (and to
    \texttt{SENTENCE} and \texttt{ENTITY} nodes via mean-pooling). This is the
    step that turns a 13-dimensional one-hot into the 781-dimensional node
    feature consumed by the GNN.
    \item The exporter writes the graph as GraphSON or PyTorch Geometric
    tensors.
\end{enumerate}

\subsection{Security Pipeline}

The security pipeline has three practical variants. All three extend, rather
than replace, the generic pipeline above: the spaCy frontend always runs
first, and the security frontends overlay security-specific entities and
relations on top of the generic linguistic graph.

\begin{table}[h]
\centering
\small
\begin{tabular}{lll}
\toprule
Pipeline & Frontend & Role \\
\midrule
\texttt{SecurityPipeline} & rule-based & Regex and dictionary extraction. \\
\texttt{ModelSecurityPipeline} & SecBERT-based & Embeddings and similarity-based extraction. \\
\texttt{HybridSecurityPipeline} & rule + model & Default security pipeline when model dependencies are available. \\
\bottomrule
\end{tabular}
\caption{Security pipeline variants in the implementation.}
\end{table}

The rule-based frontend is deterministic and strong for structured identifiers
such as CVE IDs, CWE IDs, and version strings. The model-based frontend adds
contextual embeddings and can identify less rigid security language. The hybrid
frontend runs rule extraction first and then adds model-derived entities that
do not overlap with rule-derived spans. This gives a practical balance between
precision on canonical patterns and coverage for less standard wording.

\section{TPG versus Tokenization}

Tokenization is necessary, but it is not a complete representation of text. A
tokenizer answers the question: what units does the text split into? A TPG
answers a larger question: what typed objects and typed relations are present
in the text?

\begin{table}[h]
\centering
\small
\begin{tabular}{p{0.28\textwidth}p{0.30\textwidth}p{0.30\textwidth}}
\toprule
Aspect & Tokenization & Text Property Graph \\
\midrule
Primary unit & token or subword & typed node and typed edge \\
Structure & linear order & multi-relational graph \\
Syntax & implicit or absent & explicit \texttt{DEP} edges \\
Coreference & not represented directly & explicit \texttt{COREF} edges \\
Document layout & usually flattened & document, paragraph, sentence hierarchy \\
Domain relations & not represented directly & optional security nodes and edges \\
Model input & sequence tensor & graph tensors plus node features \\
\bottomrule
\end{tabular}
\caption{Difference between tokenization and TPG representation.}
\end{table}

The TPG still uses tokens. Token nodes are part of the graph, and transformer
embeddings may be attached to those nodes. The difference is that token
embeddings are not the whole representation. They become node features inside a
larger graph that preserves structural and domain relations.

\section{TPG versus Knowledge Graphs}

A knowledge graph usually stores normalized facts about entities. For example,
it may store that one CVE affects one product or has one severity. This is
useful, but a knowledge graph often abstracts away the text that produced the
fact.

A TPG keeps the linguistic evidence and the extracted domain relations in the
same graph. The sentence, tokens, predicates, entity spans, dependency edges,
and security relation edges are present together. This makes TPG closer to a
text analysis graph than to a clean factual database.

\begin{table}[h]
\centering
\small
\begin{tabular}{p{0.28\textwidth}p{0.30\textwidth}p{0.30\textwidth}}
\toprule
Aspect & Knowledge graph & Text Property Graph \\
\midrule
Main goal & Store facts about entities & Preserve text structure and extracted relations \\
Typical node & canonical entity & document, sentence, token, entity, predicate, security object \\
Text evidence & often external or discarded & retained inside the graph \\
Syntax and discourse & usually absent & explicit edge layers \\
Best use & factual lookup and reasoning & graph learning, structural retrieval, explainable text analysis \\
\bottomrule
\end{tabular}
\caption{Knowledge graph versus TPG.}
\end{table}

The security TPG can contain knowledge-graph-like facts, but it is not only a
knowledge graph. It is a text-derived property graph with linguistic, document,
and security layers.

\section{How TPG Helps Security Text}

Security meaning is relational. A useful representation should preserve
relations such as:

\begin{itemize}
    \item CVE affects software;
    \item software has version;
    \item CVE is classified as CWE;
    \item attack vector causes impact;
    \item vulnerability is mitigated by remediation;
    \item CVE has severity;
    \item vulnerability mentions a code element.
\end{itemize}

These relations are represented directly by the security edge types. They give
a graph neural network more information than a list of security keywords. The
model can pass messages along typed edges such as \texttt{SEC\_AFFECTS},
\texttt{SEC\_EXPLOITED\_BY}, and \texttt{SEC\_MITIGATED\_BY}. This allows the
graph branch to distinguish between the presence of a term and the role that
term plays in the advisory.

The EPSS-TPG model also uses a multiview GNN. In that setup, edge types are
grouped into views such as syntactic, sequential, semantic, discourse, and
security views. This matches the CPG-inspired design: different relations are
not collapsed into one adjacency matrix, but kept as separate signals that the
model can weight.

\section{Worked Example}

Consider the sentence:

\begin{quote}
\textit{CVE-2024-1234: A buffer overflow in libfoo 1.2.3 allows remote
attackers to execute arbitrary code. The vulnerability has a CVSS v3 base
score of 9.8 indicating high severity.}
\end{quote}

After the linguistic frontend, security frontend, and security-relations pass,
the graph contains ordinary linguistic nodes and security nodes, including:

\begin{itemize}
    \item a \texttt{CVE\_ID} node for \texttt{CVE-2024-1234};
    \item a \texttt{VULN\_TYPE} node for \texttt{buffer overflow};
    \item a \texttt{SOFTWARE} node for \texttt{libfoo};
    \item a \texttt{VERSION} node for \texttt{1.2.3};
    \item an \texttt{ATTACK\_VECTOR} node for \texttt{remote};
    \item an \texttt{IMPACT} node for \texttt{execute arbitrary code};
    \item a \texttt{SEVERITY} node for \texttt{high severity}.
\end{itemize}

The graph can then contain security edges such as:

\begin{itemize}
    \item \texttt{SEC\_AFFECTS} from \texttt{CVE-2024-1234} to
    \texttt{libfoo};
    \item \texttt{SEC\_HAS\_VERSION} from \texttt{libfoo} to
    \texttt{1.2.3};
    \item \texttt{SEC\_EXPLOITED\_BY} from \texttt{CVE-2024-1234} to
    \texttt{remote};
    \item \texttt{SEC\_CAUSES} from \texttt{remote} to
    \texttt{execute arbitrary code};
    \item \texttt{SEC\_HAS\_SEVERITY} from \texttt{CVE-2024-1234} to
    \texttt{high severity}.
\end{itemize}

In addition, every \texttt{TOKEN} node carries a 768-dimensional SecBERT
embedding obtained by averaging the WordPiece sub-token vectors whose character
spans fall inside the token. The \texttt{SENTENCE} and \texttt{ENTITY} nodes
carry sentence- and entity-level embeddings derived in the same step. The
final node-feature vector consumed by the GNN is therefore the concatenation
of the 13-dimensional node-type one-hot and the 768-dimensional contextual
embedding, for a total of 781 dimensions per node.

This example shows the difference between text extraction and graph
representation. The TPG does not only detect the words; it connects them with
typed relations \emph{and} attaches contextual embeddings to those connected
nodes, so a downstream GNN can reason over both the structure and the local
semantics simultaneously.

\section{Positioning as a State-of-the-Art-Style Technique}

The TPG should be described carefully. It is not enough to say that any graph
representation is automatically state of the art. The contribution is more
specific. TPG combines three modern ideas in one implemented pipeline:

\begin{enumerate}
    \item it adopts the CPG principle of multi-view representation and
    transfers it to text;
    \item it keeps transformer embeddings as node features rather than
    discarding neural language representations;
    \item it adds a domain-specific security layer so the graph contains both
    linguistic structure and security semantics.
\end{enumerate}

This makes TPG stronger than plain tokenization, richer than a single
dependency tree, and more evidence-preserving than a clean knowledge graph. A
fair claim is that this project presents a modern, CPG-inspired, multi-view
graph representation for security text. Downstream performance claims should
still be tied to the actual experiments.

\section{Limitations}

The TPG is useful, but it should not be overstated.

\begin{itemize}
    \item Coreference is heuristic in the current enrichment pass. It uses
    recency and basic agreement rather than a full neural coreference system.
    \item Discourse relations are marker-based. They capture explicit markers
    such as \texttt{however}, \texttt{because}, and \texttt{subsequently}, but
    they do not perform full discourse parsing.
    \item Security extraction depends on rules, dictionaries, thresholds, and
    model coverage.
    \item Security edges are opt-in. If the pipeline does not enable security
    relations and the GNN edge vocabulary is not expanded, the model cannot use
    first-class \texttt{SEC\_*} edge types.
    \item TPG is a representation. It improves the information available to a
    model, but it does not guarantee better performance on every downstream
    task.
\end{itemize}

\section{Conclusion}

The Text Property Graph represents text as a structured, typed, attributed
graph. Its central idea is borrowed from the Code Property Graph: keep several
views of the input in one graph rather than treating them as separate analyses.
For text, those views are dependency structure, sequence, coreference,
predicate-argument structure, discourse, entity relations, and optional
security-specific relations.

This makes TPG different from tokenization and different from a conventional
knowledge graph. Tokenization gives a sequence of units. A knowledge graph
usually gives normalized facts. A TPG keeps both the textual evidence and the
relations extracted from that evidence.

For security-related text, this is especially useful because the meaning of an
advisory is carried by relations: which vulnerability affects which software,
through which attack vector, with which impact, and with which remediation.
The implemented TPG pipeline therefore provides a defensible, CPG-inspired,
state-of-the-art-style representation for security text.

\end{document}
